% Original analytic teaching diagram, 2026-11-12.
% Which representation exists for which kind of distribution on the real line.
\definecolor{VB}{HTML}{7188A3}\definecolor{VBF}{HTML}{E5EAF0}
\definecolor{VG}{HTML}{829C8B}\definecolor{VGF}{HTML}{E8EEE9}
\definecolor{VY}{HTML}{B4A16F}\definecolor{VYF}{HTML}{F1ECDD}
\definecolor{VR}{HTML}{A7787D}\definecolor{VRF}{HTML}{F0E5E6}
\definecolor{VP}{HTML}{9588A5}\definecolor{VPF}{HTML}{ECE8F0}
\definecolor{VI}{HTML}{3B4252}\definecolor{VM}{HTML}{6B7280}\definecolor{VL}{HTML}{4C4D4F}
\begin{tikzpicture}[x=1mm,y=1mm,>=Stealth,font=\figurefont\fontsize{8}{10}\selectfont,
lab/.style={text=VM,font=\figurefont\fontsize{7.5}{9}\selectfont}]
% outer box: all distributions (distribution function always exists)
\draw[VL,line width=.8pt,rounded corners=2pt] (0,0) rectangle (104,50);
\node[anchor=north west,text=VI] at (2,49) {全部分布：分布函数$F_X$};
% continuous F region
\draw[VY,line width=.9pt,rounded corners=2pt] (50,3) rectangle (101,40);
\node[anchor=north west,text=VI] at (52,39) {$F_X$连续};
% density region inside continuous
\fill[VBF,rounded corners=2pt] (53,14) rectangle (98,32);
\draw[VB,line width=.9pt,rounded corners=2pt] (53,14) rectangle (98,32);
\node[text=VI] at (75.5,26) {有密度$f_X$};
\node[lab] at (75.5,19.5) {等待时间、圆盘半径};
\node[lab] at (75.5,8) {Cantor分布};
% discrete region
\fill[VBF,rounded corners=2pt] (4,14) rectangle (44,32);
\draw[VB,line width=.9pt,rounded corners=2pt] (4,14) rectangle (44,32);
\node[text=VI] at (24,26) {有概率表$p_X$};
\node[lab] at (24,19.5) {骰子、抽检、首次六点};
% mixed
\node[lab] at (24,8) {等待$W$、截断$\min(T,2)$};
\end{tikzpicture}
